\documentclass[10pt,a4paper]{amsart}
\usepackage[margin=19mm]{geometry}
\usepackage{amsmath,amssymb,mathtools}
\usepackage[scr=rsfs,cal=euler]{mathalfa}
\usepackage{xfrac}
\usepackage[unicode,hidelinks,bookmarks=false]{hyperref}
\newcommand{\ie}{i.e.\ }
\newcommand{\cf}{cf.\ }
\newcommand{\resp}{respectively\ }
\newcommand{\Subsection}{\subsection}
\providecommand{\nobreakdash}{\nobreak\hbox{-}}
\setlength{\emergencystretch}{2em}
\multlinegap0pt
\usepackage{bm}
\usepackage{bbm}
\usepackage{dsfont}
\usepackage{xspace}
\usepackage{cancel}
\usepackage{mathtools}
\usepackage{amsthm}
\makeatletter
\@ifundefined{theorem}{\newtheorem{theorem}{Theorem}}{}
\@ifundefined{lemma}{\newtheorem{lemma}[theorem]{Lemma}}{}
\@ifundefined{proposition}{\newtheorem{proposition}[theorem]{Proposition}}{}
\makeatother
\newcommand{\AG}{{\mathrm{AG}}}
\newcommand{\A}{{\mathcal{A}}} %weight sequence%
\newcommand{\B}{{\mathcal{B}}} %weight sequence%
\newcommand{\Mj}{{\{M_j\}_{j\in\mathbb{N}_0}}} %weight sequence%
\newcommand{\M}{{\mathcal{M}}} %weight sequence%
\newcommand{\N}{{\mathbb{N}}} %natural numbers%
\title[Gelfand-Shilov spaces and operators with ultradifferential w]{Gelfand-Shilov spaces and operators with ultradifferential weighted symbols on non-compact manifolds}
\author{Sandro Coriasco, Pedro Meyer Tokoro}
\date{Source: arXiv:2606.27327v1 (CC BY 4.0)}
\begin{document}
\maketitle

\noindent\emph{Source and licence.} Gelfand-Shilov spaces and operators with ultradifferential weighted symbols on non-compact manifolds. Version arXiv:2606.27327v1 (CC BY 4.0), distributed under the Creative Commons Attribution 4.0 International licence, as recorded in the arXiv licence metadata for this exact version and shown on the abstract page https://arxiv.org/abs/2606.27327v1. Full rights notice: licenses/arxiv-2606.27327-CC-BY-4.0.txt.\par\smallskip

\noindent\textit{Editorial scope.} Counted unit: the Lemma whose source label is \texttt{est\_FGH} in the article (source0.tex lines 2114--2125, source label \texttt{est\_FGH}), delivered together with the article's own complete original proof of it (source0.tex lines 2127--2302, one \texttt{proof} environment of 176 lines). No mathematics is written, repaired or re-derived: statement and proof are the article's own text line for line. Required context retained uncounted: incl\_quasi\_factorial (lines 295--297). Every definition, display and label of those lines is retained unchanged. Omitted from this package and not redistributed: 1--294, 298--2113, 2126--2126, 2303--3081. Nothing inside a retained excerpt is omitted or altered: every retained body is delivered exactly as the article's lines are written, so no omission receipt of any kind accompanies this package. Citations: the retained excerpts carry no citation command at all, so no bibliography entry of the article is transcribed and no reference is redistributed. Reference adaptation: none was needed and none was made; every reference inside the delivered unit resolves inside it, so no unresolved reference or citation is produced. Printed slips kept literal: no printed word, symbol, hypothesis or number of the counted statement or of its proof was altered. Layout and class: the article's own class is \texttt{amsart} with packages including indentfirst, enumerate, cite, amssymb, amsfonts, amsmath, amsthm, mathrsfs, dsfont, color, graphicx, multicol, hyperref, geometry; the wrapper is the lane's pinned \texttt{amsart} editorial wrapper at 10pt on A4, which restates the theorem-like environments the retained text needs and adds the article's own macro definitions that the retained text needs (\texttt{\textbackslash A}, \texttt{\textbackslash AG}, \texttt{\textbackslash B}, \texttt{\textbackslash M}, \texttt{\textbackslash Mj}, \texttt{\textbackslash N}), with extra wrapper packages (bm, bbm, dsfont, xspace, cancel, mathtools). The delivered numbering is the unit's own. Assets: no retained span contains a figure environment, a graphics-inclusion command, a PDF-page inclusion, a table or a TikZ picture; no image file is copied into this package, the package declares no asset dependency and no image file is redistributed with it. Licence: version arXiv:2606.27327v1 of this article is distributed under the Creative Commons Attribution 4.0 International licence, as recorded in the arXiv licence metadata for this exact version and shown on the abstract page https://arxiv.org/abs/2606.27327v1. A full rights notice is delivered at licenses/arxiv-2606.27327-CC-BY-4.0.txt. Characters: every retained excerpt is pure ASCII, so the delivered engine, XeLaTeX with its default Latin Modern text font, reports no missing character.
\begin{proposition}\label{incl_quasi_factorial}
	Suppose that $\M=\Mj$ satisfies $(M3)'$. Then $\{j!\}_{j\in\N_0}\prec \M$.
\end{proposition}

\begin{lemma}\label{est_FGH}
	For every fixed $t\in[0,1]$, there exist $C,h>0$ 
	(respectively, for every $h>0$, there exists $C>0$) such that, for all $\alpha,\beta,\gamma\in\mathbb{N}_0^n$ and $x,y,\xi\in\mathbb{R}^n$, we have:
	%
	\begin{enumerate}
		\item[(a)] $|\partial_x^\beta\partial_y^\gamma G_t(x,y)|\leq Ch^{|\beta|+|\gamma|}B_{|\beta|}B_{|\gamma|}\langle x\rangle^{-|\beta|-|\gamma|}\langle x-y\rangle^{|\beta|+|\gamma|}$;
		
		\item[(b)] $|\partial_x^\beta\partial_y^\gamma H_t(x,y)|\leq Ch^{|\beta|+|\gamma|}B_{|\beta|}B_{|\gamma|}\langle x\rangle^{-|\beta|-|\gamma|}\langle x-y\rangle^{|\beta|+|\gamma|}$;
		
		\item[(c)] $|\partial_\xi^\alpha\partial_x^\beta\partial_y^\gamma F_t(x,y,\xi)|\leq Ch^{|\alpha|+|\beta|+|\gamma|}A_{|\alpha|}B_{|\beta|}B_{|\gamma|} \langle\xi\rangle^{m_1-|\alpha|} \langle x\rangle^{m_2-|\beta|-|\gamma|}\langle x-y\rangle^{|\beta|+|\gamma|}$.
	\end{enumerate}
\end{lemma}

\begin{proof}
	First, by Peetre's inequality, we have that
	%
	\begin{equation}\label{peetre}
		\langle x-t(y-x)\rangle^{-j} \leq 2^j\langle x\rangle^{-j}\langle x-y\rangle^{j},
	\end{equation}
	%
	for all $t\in[0,1]$, $x,y\in\mathbb{R}^n$, and $j\geq 0$.
	
	\medskip\noindent
	(a) Write
	\[\varphi' = \begin{bmatrix} \partial_{x_1}\varphi_1 & \cdots & \partial_{x_n}\varphi_1 \\ \vdots & \ddots & \vdots \\ \partial_{x_1}\varphi_n & \cdots & \partial_{x_n}\varphi_n \end{bmatrix} = \begin{bmatrix} \varphi_1^{(1)} & \cdots & \varphi_1^{(n)} \\ \vdots & \ddots & \vdots \\ \varphi_n^{(1)} & \cdots & \varphi_n^{(n)} \end{bmatrix}.\]
	
	Then, we have
	\[\det\varphi' = \sum_{\sigma\in S_n}\mathrm{sgn}(\sigma)\varphi_1^{(\sigma(1))}\cdots\varphi_n^{(\sigma(n))},\]
	where $S_n$ is the symmetric group of order $n$, which has $n!$ elements. By $(I1)$, there exist $C,h>0$ such that
	\[|\partial_x^\alpha\varphi_\ell^{(\sigma(\ell))}(x)|\leq Ch^{|\alpha|}\alpha!\langle x\rangle^{-|\alpha|},\]
	for all $\alpha\in\mathbb{N}_0^n$, $x\in\mathbb{R}^n$, $\ell=1,\dots,n$, and $\sigma\in S_n$. Let $\alpha_1,\cdots,\alpha_n\in\mathbb{N}_0^n$ be such that $\alpha_1+\cdots+\alpha_n=\alpha$. Then, we have
	\begin{equation*}
		|\partial_x^{\alpha_1}\varphi_1^{(\sigma(1))}\cdots \partial_x^{\alpha_n}\varphi_n^{(\sigma(n))}| \leq (Ch^{|\alpha_1|}\alpha_1!\langle x\rangle^{-|\alpha_1|})\cdots (Ch^{|\alpha_n|}\alpha_n!\langle x\rangle^{-|\alpha_n|}) \leq C^nh^{|\alpha|}\alpha!\langle x\rangle^{-|\alpha|}.
	\end{equation*}
	
	Given $j,k\in\{1,\dots,n\}$, we also see that
	\[\partial_{x_j}\det\varphi'(x) = 
	\sum_{\sigma\in S_n}\mathrm{sgn}(\sigma)\sum_{\ell=1}^n \varphi_1^{(\sigma(1))}\cdots \partial_{x_j}\varphi_{\ell}^{(\sigma(\ell))}\cdots \varphi_n^{(\sigma(n))},\]
	and
	\[\partial_{x_k}\partial_{x_j}\det\varphi'(x) = \sum_{\sigma\in S_n}\mathrm{sgn}(\sigma)\sum_{\nu=1}^n \sum_{\ell=1}^n \varphi_1^{(\sigma(1))}\cdots \partial_{x_k}\varphi_{\nu}^{(\sigma(\nu))} \cdots \partial_{x_j}\varphi_{\ell}^{(\sigma(\ell))}\cdots \varphi_n^{(\sigma(n))}.\]
	
	By induction, it follows that $\partial_x^\alpha\det\varphi'(x)$ is a linear combination of $n!n^{|\alpha|}$ terms of the form
	\[\mathrm{sgn}(\sigma)\partial_x^{\alpha_1}\varphi_1^{(\sigma(1))}\cdots \partial_x^{\alpha_n}\varphi_n^{(\sigma(n))},\]
	with $\alpha_1+\cdots+\alpha_n=\alpha$. Hence, we conclude
	\[|\partial_x^\alpha\det\varphi'(x)|\leq C_0h_0^{|\alpha|}\alpha!\langle x\rangle^{-|\alpha|},\]
	for all $\alpha\in\mathbb{N}_0^n$ and $x\in\mathbb{R}^n$, where $C_0=C^nn!$ and $h_0=nh$.
	From this estimate and \eqref{peetre}, we obtain $C,h>0$ such that
	%
	\[|\partial_x^\beta\partial_y^\gamma G_t(x,y)|\leq Ch^{|\beta|+|\gamma|}(\beta!\gamma!)\langle x\rangle^{-|\beta|-|\gamma|}\langle x-y\rangle^{|\beta|+|\gamma|},\]
	%
	for all $\beta,\gamma\in\mathbb{N}_0^n$, $x,y\in\mathbb{R}^n$, and $t\in [0,1]$. Notice that the absolute value in the definition of $G_t(x,y)$ does not affect the smoothness of the map since $\varphi$ is a diffeomorphism. Finally, we apply Proposition \ref{incl_quasi_factorial} to obtain the desired estimate.
	
	\medskip\noindent
	(b) Using the notation $\varphi_{\ell}^{(k)}=\partial_{x_k}\varphi_\ell$ from the previous point, by $(I1)$ and \eqref{peetre} it follows
	\begin{align*}
		\left|\partial_x^\beta\int_0^1\varphi_{\ell}^{(\sigma(\ell))}(x+t(y-x))\,\mathrm{d}t\right| & 
		\leq \int_0^1\left|\partial_x^\beta \varphi_{\ell}^{(\sigma(\ell))}(x+t(y-x))\right|\,\mathrm{d}t \\
		& \leq C(2h)^{|\beta|}\beta!\langle x\rangle^{-|\beta|}\langle x-y\rangle^{|\beta|}.
	\end{align*}
	
	For the $y$-derivatives, we obtain analogous estimates. Then, the estimates for $H_t(x,y)$ follow from those obtained in (a) 
	and Proposition \ref{incl_quasi_factorial}.
	
	\medskip\noindent
	(c) Fix $t\in [0,1]$ and write $X=(X_1,\dots,X_n)$, $Y=(Y_1,\dots,Y_n)$, $\Xi=(\Xi_1,\dots,\Xi_n)$, where
	%
	\[X_j=\varphi_j(x),\quad Y_j=\varphi_j(x+t(y-x)),\quad\text{and}\quad \Xi_j=\sum_{k=1}^n \nabla\varphi(x,y)_{jk}^{-T}\xi_j,\quad j=1,\dots,n.\]
	
	In the sequel, we write
	%
	\[
		p_1^\sharp(x,y,\xi)=p_1(\varphi(x),\varphi(x+t(y-x)),\nabla\varphi(x,y)^{-T}\xi)=p_1(X,Y,\Xi).
	\]
	
	By the chain rule, 
	%
	\begin{equation}\label{dxp}
	\begin{aligned}
		\partial_{x_j}p_1^\sharp &= \sum_{k=1}^n \frac{\partial p_1}{\partial X_k}\frac{\partial X_k}{\partial x_j} + 
		\sum_{k=1}^{n}\frac{\partial p_1}{\partial Y_k}\frac{\partial Y_k}{\partial x_j} + 
		\sum_{k=1}^{n}\frac{\partial p_1}{\partial \Xi_k}\frac{\partial \Xi_k}{\partial x_j},
	\\
		\partial_{y_j}p_1^\sharp &= \sum_{k=1}^{n}\frac{\partial p_1}{\partial Y_k}\frac{\partial Y_k}{\partial y_j} + \sum_{k=1}^{n}\frac{\partial p_1}{\partial \Xi_k}\frac{\partial \Xi_k}{\partial y_j},
	\\
		\partial_{\xi_j}p_1^\sharp &= \sum_{k=1}^n \frac{\partial p_1}{\partial\Xi_k}\nabla\varphi(x,y)_{kj}^{-T}.
	\end{aligned}
	\end{equation}
	
	Observe that, by $(I1)$, it holds
	\[
		\langle\varphi(x)\rangle \asymp \langle x\rangle\quad\text{and}\quad \langle \nabla\varphi(x,y)^{-T}\xi\rangle \asymp \langle\xi\rangle.
	\]
	
	Then, the fact that $p_1\in \AG_{[\A,\B]}^{m,0,\mu}(\mathbb{R}^n)$ and \eqref{peetre}, together with $(I1)$, imply
	%
	\begin{equation*}
		\left|\frac{\partial p_1}{\partial X_j}\right| \leq Ch B_1\langle x\rangle^{m-1}  \langle \xi\rangle^{\mu},\quad 
		\left|\frac{\partial p_1}{\partial Y_j}\right| \leq Ch B_1 \langle x\rangle^{m-1} \langle x-y\rangle \langle \xi\rangle^{\mu},
	\end{equation*}
	%
	and
	%
	\begin{equation*}
		\left|\frac{\partial p_1}{\partial \Xi_j}\right| \leq Ch A_1 \langle x\rangle^{m} \langle\xi\rangle^{\mu-1}.
	\end{equation*}
	
	We also have, again by $(I1)$,
	%
	\begin{equation*}
		\left|\frac{\partial X_j}{\partial x_j}\right| \leq Ch B_1 ,\quad %\langle x\rangle^{1-1},\quad 
		\left|\frac{\partial Y_j}{\partial x_j}\right| \leq Ch B_1, %\langle x\rangle^{1-1}\langle x-y\rangle^{1-1},
	\end{equation*}
	%
	\begin{equation*}
		\left| \frac{\nabla\varphi(x,y)_{k\ell}^{-T}}{\partial x_j}\right| \leq Ch B_1\langle x\rangle^{-1}\langle x-y\rangle,
		\quad \left|\frac{\partial Y_j}{\partial y_j}\right| \leq Ch B_1%\langle x\rangle^{1-1}\langle x-y\rangle^{1-1},
	\end{equation*}
	%
	\begin{equation*}
		\left|\frac{\nabla\varphi(x,y)_{k\ell}^{-T}}{\partial y_j}\right| \leq Ch B_1\langle x\rangle^{-1}\langle x-y\rangle,\quad |\nabla\varphi(x,y)_{kj}^{-T}|\leq C.
	\end{equation*}
	
	Notice that the terms $\xi_\ell$ that arise when we differentiate $p_1^\sharp$ with respect to $x_j$ or $y_j$ are compensated by the decay given by 
	$\partial p_1^\sharp/\partial\Xi_\ell$. Then, combining the previous estimates, we have proven
	%
	\begin{equation}\label{est_1st}
		\left| \partial^\alpha_x\partial^\beta_y\partial^\gamma_\xi p_1^\sharp(x,y,\xi) \right|\leq (C^23nh)h^1 A_{|\alpha|}B_{|\beta|}B_{|\gamma|} 
		 \langle x\rangle^{m-|\alpha|-|\beta|}\langle x-y\rangle^{|\alpha|+|\beta|} \langle\xi\rangle^{\mu-|\gamma|},
	\end{equation}
	%
	for $|\alpha|+|\beta|+|\gamma|\leq 1$.
	
	If we differentiate \eqref{dxp} with respect to $x_\ell$, the Leibniz rule yields
	%
	\begin{align*}\label{dxxp}
		\partial_{x_\ell}\partial_{x_j}p_1^\sharp = & \sum_{k=1}^n \left( \left(\partial_{x_\ell}\frac{\partial p_1}{\partial X_k}\right)\frac{\partial X_k}{\partial x_j}+\frac{\partial p_1}{\partial X_k}\frac{\partial^2 X_k}{\partial x_\ell\partial x_j}\right) + \sum_{k=1}^{n}\left( \left(\partial_{x_\ell}\frac{\partial p_1}{\partial Y_k}\right)\frac{\partial Y_k}{\partial x_j}+\frac{\partial p_1}{\partial Y_k}\frac{\partial^2 Y_k}{\partial x_\ell\partial x_j}\right)\\
		 + &\sum_{k=1}^{n}\left( \left(\partial_{x_\ell}\frac{\partial p_1}{\partial \Xi_k}\right)\frac{\partial \Xi_k}{\partial x_j}+\frac{\partial p_1}{\partial \Xi_k}\frac{\partial^2 \Xi_k}{\partial x_\ell\partial x_j}\right). 
	\end{align*}
	
	Since $p_1\in \AG_{[\A,\B]}^{m,0,\mu}$, it follows
	%
	\[
	\frac{\partial p_1}{\partial X_k}(X,Y,\Xi)\in \AG_{[\A,\B]}^{m-1,0,\mu},\ 
	\frac{\partial p_1}{\partial Y_k}(X,Y,\Xi)\in \AG_{[\A,\B]}^{m,-1,\mu},\ 
	\frac{\partial p_1}{\partial \Xi_k}(X,Y,\Xi)\in \AG_{[\A,\B]}^{m,0,\mu-1}.
	\]
	
	Moreover, these symbols satisfy
	%
	\begin{equation*}
		\left|\partial^\alpha_X\partial^\beta_Y\partial^\gamma_\Xi
		\frac{\partial p_1}{\partial X_k}\right|
		\leq Ch^{|\alpha|+|\beta|+|\gamma|+1} A_{|\alpha|}B_{|\beta|+1}B_{|\gamma|} 
		\langle x\rangle^{m-|\alpha|-1}\langle y\rangle^{-|\beta|} \langle \xi\rangle^{\mu-|\gamma|},
	\end{equation*}
	%
	\begin{equation*}
		\left|\partial^\alpha_X\partial^\beta_Y\partial^\gamma_\Xi
		\frac{\partial p_1}{\partial Y_k}\right|
		\leq Ch^{|\alpha|+|\beta|+|\gamma|+1}A_{|\alpha|} B_{|\beta|}B_{|\gamma|+1} 
		\langle x\rangle^{m-|\alpha|}\langle y\rangle^{-|\beta|-1}\langle \xi\rangle^{\mu-|\gamma|},
	\end{equation*}
	%
	\begin{equation*}
		\left|
		\partial^\alpha_X\partial^\beta_Y\partial^\gamma_\Xi
		\frac{\partial p_1}{\partial \Xi_k}\right|\leq Ch^{|\alpha|+|\beta|+|\gamma|+1}A_{|\alpha|+1} B_{|\beta|}B_{|\gamma|} 
		\langle x\rangle^{m-|\alpha|}\langle y\rangle^{-|\beta|}\langle \xi\rangle^{\mu-|\gamma|-1},
	\end{equation*}
	%
	for $\alpha,\beta,\gamma\in\mathbb{N}_0^n$, where the constants $C,h>0$ are the same appearing in the previous estimates above.
	If we now apply \eqref{est_1st} to those symbols in place of $p_1^\sharp$, we obtain
	%
	\begin{align*}
		\left| \partial_{x_\ell}\frac{\partial p_1}{\partial X_k} \right| &\leq (C^23nh)h^2  B_2 \langle x\rangle^{m-2}\langle \xi\rangle^{\mu},
	\\		
		\left| \partial_{x_\ell}\frac{\partial p_1}{\partial Y_k} \right| &\leq (C^23nh)h^2 B_2\langle x\rangle^{m-2}\langle x-y\rangle^2 \langle \xi\rangle^{\mu},
	\\
		\left| \partial_{x_\ell}\frac{\partial p_1}{\partial \Xi_k} \right| &\leq (C^23nh)h^2 A_1 B_1 \langle x\rangle^{m-1}\langle x-y\rangle^2.
	\end{align*}
	
	By $(I1)$ for the derivatives of $X_j$, $Y_j$, and $\Xi_j$, $\partial_{x_\ell}\partial_{x_j}p_1^\sharp$ is a sum of $2(3n)$ terms satisfying the desired estimates. 
	By an induction argument, a derivative of order $N$ is a sum of $2^{N-1}3n$ terms. Then, incorporating the factor $1/2$ within the constant $C_0=C^23nh$ and 
	$2^N$ within the constant $h$, we obtain the desired estimate for the $x$-derivatives. Similar arguments give the estimates for the derivatives with respect to 
	$y$ and $\xi$. This proves the estimates in the inductive case.
	
	The estimates for the projective case can be proved in a similar fashion, by choosing the constants 
	$h>0$ in an appropriate way, to absorb the extra constants that appear during the computations. The proof is complete.
\end{proof}

\end{document}
